\subsubsection{Formal Statement}\label{sec:no-cloning-formal-statement}

\begin{definitionbox}[: No-Cloning Theorem]
    The \definition{No-Cloning Theorem} states that there \textbf{does not exist a universal quantum gate capable of perfectly copying an arbitrary unknown quantum state}.

    \highspace
    Formally, suppose we have:
    \begin{itemize}
        \item A qubit in an unknown state: \ket{x}.
        \item And a second qubit initialized to: \ket{0}.
    \end{itemize}
    We would like a quantum gate $U$ such that:
    \begin{equation}
        U \ket{x} \ket{0} = \ket{x} \ket{x}
    \end{equation}
    For every possible quantum state \ket{x}. The \hl{No-Cloning Theorem states that such a gate $U$ cannot exist.}
\end{definitionbox}

\begin{proof}
    We now prove that the cloning gate $U$:
    \begin{equation*}
        U \ket{x} \ket{0} = \ket{x} \ket{x}
    \end{equation*}
    Cannot exist for an arbitrary quantum state \ket{x}.
    \begin{enumerate}
        \item \textbf{Assume that cloning is possible}. Suppose there exists a unitary gate $U$ such that:
        \begin{equation*}
            U \ket{x} \ket{0} = \ket{x} \ket{x}
        \end{equation*}
        And also:
        \begin{equation*}
            U \ket{y} \ket{0} = \ket{y} \ket{y}
        \end{equation*}
        So $U$ should clone both \ket{x} and \ket{y}.

        \item \textbf{Take the inner product before cloning}. Before cloning, the two input states are:
        \begin{equation*}
            \ket{x}\ket{0} \quad \text{and} \quad \ket{y}\ket{0}
        \end{equation*}
        Their inner product is:
        \begin{equation*}
            \braket{x}{y} \braket{0}{0}
        \end{equation*}
        Since \ket{0} is a normalized state, we have $\braket{0}{0} = 1$. Thus, the inner product before cloning is:
        \begin{equation*}
            \braket{x}{y}
        \end{equation*}

        \item \textbf{Take the inner product after cloning}. After cloning, the states become:
        \begin{equation*}
            \ket{x}\ket{x} \quad \text{and} \quad \ket{y}\ket{y}
        \end{equation*}
        Their inner product is:
        \begin{equation*}
            \braket{x}{y} \braket{x}{y} = \left(\braket{x}{y}\right)^2
        \end{equation*}
        Now we use the tensor product property of inner products, which states that (\autopageref{eq:inner-product-tensor}):
        \begin{equation*}
            \left(A \otimes B\right)\left(C \otimes D\right) = \left(AC\right) \otimes \left(BD\right)
        \end{equation*}
        Applying this to our inner product, we get:
        \begin{equation*}
            \left(\bra{x} \otimes \bra{x}\right) \left(\ket{y} \otimes \ket{y}\right) = \left(\braket{x}{y}\right) \otimes \left(\braket{x}{y}\right)
        \end{equation*}
        Which simplifies to:
        \begin{equation*}
            \left(\braket{x}{y}\right)^2
        \end{equation*}

        \item \textbf{Use unitary}. Valid quantum gates must be unitary, which means they preserve inner products. Therefore, the inner product before and after applying $U$ must be the same:
        \begin{equation*}
            \braket{x}{y} = \left(\braket{x}{y}\right)^2
        \end{equation*}
        This is only possible if $\braket{x}{y} = 0$ or $\braket{x}{y} = 1$. But for arbitrary states, this is not generally true. Thus, the assumption that cloning is possible leads to a contradiction.
    \end{enumerate}
    Let's try to understand why this is a contradiction. Suppose two states are \ket{x} and \ket{y}, and suppose $\left|\braket{x}{y}\right|$ (i.e., the magnitude of the inner product) is $0.8$. It means that the states are not identical, but they are still ``quite similar''. They are partially distinguishable, but not perfectly. Now, if cloning were possible, the inner product after cloning would be still $0.8$. However, we found that the inner product after cloning would be $\left(0.8\right)^2 = 0.64$. This means that the states have become more distinguishable after cloning, which is a contradiction because unitary operations cannot change the inner product between states:
    \begin{equation*}
        \left|\braket{x}{y}\right| = 0.8 \xrightarrow{\text{clone}} \left|\braket{x}{y}\right|^{2} = 0.64 \xrightarrow{\text{clone}} \left|\braket{x}{y}\right|^{3} = 0.512 \dots
    \end{equation*}
\end{proof}

\begin{deepeningbox}[: Alternative Proof of No-Cloning Theorem]\phantomsection\label{deep-box:alternative-proof-no-cloning}
    We can also prove the no-cloning theorem using a different approach, by directly applying the unitarity of $U$ to the cloning operation. This alternative proof is often more concise and elegant.
    
    \begin{proof}
        Suppose, by contradiction, that there exists a cloning gate $U$ such that:
        \begin{equation*}
            U\ket{x}\ket{0}=\ket{x}\ket{x}
        \end{equation*}
        For an arbitrary state $\ket{x}$.

        \highspace
        Since $U$ is supposed to be universal, it must also clone another arbitrary state $\ket{y}$:
        \begin{equation*}
            U\ket{y}\ket{0}=\ket{y}\ket{y}
        \end{equation*}
        At this point, we have two equations describing the action of $U$ on two different input states.

        \highspace
        Now, the \textbf{only thing we know about unitary operators is that they preserve inner products} (\autopageref{def:unitary-gate}). So if we want to derive a \hl{contradiction from unitarity}, we must compute an inner product. If the inner product before cloning is different from the inner product after cloning, then we have a contradiction, which means that $U$ cannot be unitary, and therefore cannot be a valid quantum gate.

        \highspace
        Let's compute the inner product before and after cloning.
        \begin{itemize}
            \item \textbf{Before cloning}, the inner product between the two input states is:
            \begin{equation*}
                \left(U\ket{x}\ket{0}\right) \cdot \left(U\ket{y}\ket{0}\right)
            \end{equation*}
            But the inner product between two states is written as (\autopageref{eq:inner-product-2}):
            \begin{align*}
                \left(U\ket{x}\ket{0}\right) \cdot \left(U\ket{y}\ket{0}\right) &= \left(U\ket{x}\ket{0}\right)^{\dagger} \left(U\ket{y}\ket{0}\right) \\[.2em]
                %
                &=
                \left(\bra{0}\bra{x}U^\dagger\right)\left(U\ket{y}\ket{0}\right) \\[.2em]
                %
                &\downarrow \text{Remove parentheses} \\[.2em]
                %
                &=
                \bra{0}\bra{x}U^\dagger U\ket{y}\ket{0} \\[.2em]
                %
                &\downarrow \text{Use unitarity: } U^\dagger U = I \\[.2em]
                %
                &=
                \bra{0}\braket{x}{y}\ket{0}
            \end{align*}

            \item \textbf{After cloning}, the inner product between the two output states is:
            \begin{equation*}
                \left(\ket{x}\ket{x}\right) \cdot \left(\ket{y}\ket{y}\right)
            \end{equation*}
            Using the tensor product property of inner products:
            \begin{align*}
                \left(\ket{x}\ket{x}\right) \cdot \left(\ket{y}\ket{y}\right) &=
                \left(\ket{x}\ket{x}\right)^{\dagger}\left(\ket{y}\ket{y}\right) \\[.2em]
                %
                &=
                \left(\bra{x}\bra{x}\right)\left(\ket{y}\ket{y}\right) \\[.2em]
                %
                &=
                \braket{x}{y} \braket{x}{y} \\[.2em]
                %
                &=
                \left(\braket{x}{y}\right)^2
            \end{align*}
        \end{itemize}
        Thus, \textbf{if cloning were possible}, we would have:
        \begin{equation*}
            \bra{0}\braket{x}{y}\ket{0} = \left(\braket{x}{y}\right)^2
        \end{equation*}
        Using the tensor product property of inner products again:
        \begin{equation*}
            \bra{\psi} \braket{\phi}{\chi} \ket{\psi} = \braket{\phi}{\chi} \braket{\psi}{\psi}
        \end{equation*}
        Remarkably, the inner product $\braket{\phi}{\chi}$ is just a complex number (\autopageref{eq:inner-product-2}), so we can pull it out of the inner product. So, we can simplify the left side:
        \begin{equation*}
            \underbrace{\bra{0}\braket{x}{y}\ket{0}}_{= \braket{x}{y} \braket{0}{0}} = \left(\braket{x}{y}\right)^{2}
        \end{equation*}
        Since $\braket{0}{0} = 1$, we have:
        \begin{equation*}
            \braket{x}{y} = \left(\braket{x}{y}\right)^{2}
            \quad
            \iff
            \quad
            \begin{cases}
                \braket{x}{y} = 0 \\
                \braket{x}{y} = 1
            \end{cases}
        \end{equation*}
        So unitarity forces the inner product to satisfy this equation. However, this equation is not true for all possible values of $\braket{x}{y}$. It only holds if $\braket{x}{y} = 0$ or $\braket{x}{y} = 1$. But for arbitrary states, $\braket{x}{y}$ can take any value between $0$ and $1$, so this is a contradiction. Therefore, a universal cloning gate $U$ cannot exist.
    \end{proof}
\end{deepeningbox}

\highspace
\begin{flushleft}
    \textcolor{Green3}{\faIcon{question-circle} \textbf{Why CNOT Does NOT Clone Qubits}}
\end{flushleft}
A common confusion is that the \hl{CNOT gate seems to copy the control qubit into the target qubit} (\autopageref{sec:controlled-not-cnot-gate}).
This is true only when the control qubit is a computational basis state, i.e. $\ket{0}$ or $\ket{1}$.
It is not true for a generic superposition.

\highspace
Classically, CNOT is often written as:
\begin{equation*}
    \text{CNOT}\left(v_A,v_B\right) = \left(v_A, v_A \oplus v_B\right)
\end{equation*}
Where $\oplus$ denotes XOR. \textcolor{Red2}{\faIcon{exclamation-triangle}} So, if the second qubit is initialized to $\ket{0}$, one \emph{may} be tempted to write:
\begin{equation*}
    \text{CNOT}\left(v_A,\ket{0}\right) = \left(v_A, v_A \oplus 0\right) = \left(v_A, v_A\right)
\end{equation*}
This is true because the XOR of any bit with $0$ is the bit itself:

\begin{table}[!htp]
    \centering
    \begin{tabular}{@{} c c c @{}}
        \toprule
        $v_A$ & $\ket{0}$ & $v_A \oplus 0$ \\
        \midrule
        $\ket{0}$ & $\ket{0}$ & $\ket{0}$ \\
        $\ket{1}$ & $\ket{0}$ & $\ket{1}$ \\
        \bottomrule
    \end{tabular}
\end{table}

\noindent
For example, if we start with $v_A = \ket{1}$ and the second qubit is $\ket{0}$, applying CNOT gives:
\begin{equation*}
    \text{CNOT}\left(\ket{1},\ket{0}\right) = \left(\ket{1}, \ket{1} \oplus \ket{0}\right) = \left(\ket{1}, \ket{1}\right)
\end{equation*}
This looks \textbf{like cloning}: the first qubit \hl{$v_A$ appears to be copied into the second qubit.}

\highspace
However, the mistake is that \hl{this notation is valid only for computational basis states $\ket{0}$ and $\ket{1}$}, not for arbitrary superpositions. The informal notation:
\begin{equation*}
    \text{CNOT}\left(v_A,v_B\right) = \left(v_A, v_A \oplus v_B\right)
\end{equation*}
Can be used only when the inputs are basis states. For a generic superposition, we must use the full linear algebra representation of CNOT, which does not clone the state.

\highspace
\textcolor{Green3}{\faIcon{check-circle} \textbf{Correct Quantum Computation.}} Let the control qubit be a generic superposition:
\begin{equation*}
    v_A = a_0\ket{0}+a_1\ket{1}
\end{equation*}
And let the target qubit be initialized to $\ket{0}$. The joint input state is:
\begin{equation*}
    v_A\ket{0}
    =
    \left( a_0\ket{0} + a_1\ket{1} \right) \ket{0}
\end{equation*}
Using the tensor product:
\begin{equation*}
    v_A\ket{0}
    =
    a_0\ket{00} + a_1\ket{10}
\end{equation*}
Now apply CNOT linearly. CNOT leaves $\ket{00}$ unchanged and maps $\ket{10}$ to $\ket{11}$:
\begin{equation*}
    \text{CNOT}\left(a_0\ket{00}+a_1\ket{10}\right)
    =
    a_0\ket{00}+a_1\ket{11}
\end{equation*}
So the real output is:
\begin{equation*}
    a_0\ket{00}+a_1\ket{11}
\end{equation*}
\hl{If CNOT had truly cloned the qubit, the output should instead be:}
\begin{equation*}
    v_Av_A
    =
    (a_0\ket{0}+a_1\ket{1})(a_0\ket{0}+a_1\ket{1})
\end{equation*}
Expanding:
\begin{equation*}
    v_Av_A
    =
    a_0^2\ket{00}
    +
    a_0a_1\ket{01}
    +
    a_1a_0\ket{10}
    +
    a_1^2\ket{11}    
\end{equation*}
However, these two states are not the same:
\begin{equation*}
    a_0\ket{00}+a_1\ket{11}
    \neq
    a_0^2\ket{00}
    +
    a_0a_1\ket{01}
    +
    a_1a_0\ket{10}
    +
    a_1^2\ket{11}
\end{equation*}
Therefore, CNOT does not clone a superposition. More fundamentally, a true universal cloning operation cannot exist because it would violate the unitary structure of quantum mechanics. Indeed, in general:
\begin{equation*}
    a_0\ket{00}+a_1\ket{11}
\end{equation*}
Cannot be factorized as a tensor product of two independent single-qubit states. The second state $v_Av_A$ would be two independent qubits in the same state. This is the key distinction: CNOT can create correlation or entanglement, but it does not create an independent copy of an unknown qubit. In summary, the \hl{CNOT output is entangled}, while the \hl{true cloned state would be two qubits in the same state.}

\highspace
In other words, CNOT preserves quantum coherence between basis components, while true cloning would require duplicating the amplitudes themselves, which is forbidden by unitary quantum evolution.